\documentclass[12pt,a4paper]{article}

\title{Volatility Forecasting: A Statistical and Deep Learning Approach}
\author{}
\date{}

\begin{document}

\maketitle

\begin{abstract}
% Maximum one page. Write this section last.
\end{abstract}


\section{Introduction}
% Write this section last.

\subsection{Context}
% Briefly describe the context of the work.
% Example: All human beings need food to survive.

Stock markets originated at 1602, in Amsterdam, Dutch Republic (now the Netherlands). In that year, the Dutch East India Company issues transferable shares, and investors began trading them with one another. At that point in time, it originated the possibility to earn money by buying the paper at a lower price, hoping that the company would grow with its expansion. Or, if someone had privileged information, e.g. a failed expedition or a sucessfull one, would be in a better position than the others, being able to sell its position before it loses value or buy more before it gets a higher price. 

The desire to obtain and interpret information that may affect stock prices remains central to financial markets. Services such as the Bloomberg Terminal give investors rapid access to news and data, while some traders use direct exchange feeds when even small delays matter. If you receive the information before every player, you have a leverage over them and can position yourself in a benefitial way. Jim Simons, the founder of Renaissance Technologies, whose fund Medallion had an average annual return of 66\% before investor fees, started his trading using news in early human trading decisions, before passing to models who studies historical market data as stated by Gregory Zuckerman in his Jim Simons biography, "The Man Who Solved the Market". 

As competition for information increases, the opportunity to profit from that same information decreases. This process leads to an efficient market. The efficient market theory was developed by Eugene Fama, which states: "Prices reflect available information". An efficient market is one where prices adjust quickly enough to available information such that investor cannot reliably earn extra returns from it after accounting for risk and trading costs. As markets converges to an extremely high degree of efficiency, one can assume that analyzing the prices historical data provides the answer to profit. As described by Gregory Zuckerman, Simons was searching historical price data for a recurring pattern. 

A historical data is known as a time series. Time series analysis has been a widely explored field in statistics and ML (Machine Learning). Statistical analysis is focused on understanding the properties of the data (stationarity, mean reversion, seasonality, trend, cycles, etc.) and creating a model, based on such characteristics that sucessfully explains the movements observed in such data. ML, in the meanwhile, is interested in creating a model that efficiently explains the data without effectively needing its traits. 

\subsubsection{Stastical Models}
An early landmark of time series statistical modelling is with Yule's work in 1927 where he introduced the AR (Autoregressive) model, where current value depends on its past values. In 1927, Slutsky demonstrated that summing random values can produce smooth, apparently cyclical movements, and in 1938, Wold extended this by showing that a stationary (constant mean, constant variance and covariance depending only on time lag) time series can be split into a predictable component and another represented by current and past shocks. Both contributions served as the foundation for the MA (Moving average) model. 

In 1970, Box and Jenkins combined AR and MA terms with differencing to model a series, originating the ARIMA (Autoregressive Integrated Moving Average) model. In ARIMA(p, d, q), p and q corresponds to the order of AR and MA, respectively, while d corresponds to the number of times the series is differenced in order to remove certain kinds of non-stationarity, especially, the one caused by a stochastic trend. In the same year, the authors extended the model to account for seasonal patterns, SARIMA (Seasonal Autoregressive Integrated Moving Average) model. SARIMA(p, d, q)(P, D, Q)s where (p, d, q) corresponds to the base ARIMA model. Seasonality is a pattern that tends to recur every s observation, e.g in a monthly data with a yearly pattern, s = 12. P and Q represent the order of the seasonal AR and MA parts, respectively, where D is the number of seasonal differences. 

Box and Jenkins's models addresses the serie's expected value over time. The unexplained data is treated as White Noise ($\varepsilon_t \sim \mathrm{WN}(0, \sigma^2)$), assuming homosckedasticity (the variance of the errors remains constant across observations). Robert Engle, twelve years later, in 1982, was motivated to measure uncertainty across the time. He was interested in studying the uncertainty about future inflation and how it impacted in the financial market. Engle was interested in studying the conditional variance, how the variance changes based on the past information present. It can be that a shock is unconditionally homosckedastic, but conditionally heterosckedastic. In order to calculate the conditional variance, Engel estimated the ARCH (Autoregressive Conditional Heterosckedastic). Autoregressive since the current variance depends on past squared shocks, in order to capture the size of a shock, regardless of direction as well. Conditional, the variance is forecast based on information available before the forecast. Heterosckedastic, the variance changes over time. 

Engel's ARCH model was revolutionary, but had a problem. Volatility often stays elevated for many periods and in order to ARCH address this behaviour it was needed many past squared schocks and consequently many parameters. The uncertainty of a market is called volatility, which represents the standard deviation of variance, volatility = $\sqrt{\sigma^2}$. Tim Bollerslev, in 1986, addresseed this limitation by creating GARCH (Generalized Autoregressive Conditional Heterosckedastic), where he added past variance values in order to let the shock effect persist without adding a long list of past squared shocks. 

In 2009, Fulvio Corsi presented the HAR (Heterogeneous Autoregressive) model. His idea was simple yet elegant: market participants act over different time horizons, so volatility measured over a day, a week, and a month may help to predict future volatility. Corsi based his idea over the Heteregeneous Market hypothesis, attributed to Ulrich A. Müller et al. (1993) in his work Fractals and Intrinsic Time: A Challenge to Econometricians. 

AR, MA, ARIMA and SARIMA models the conditional mean. HAR presents a Realized-volatility regression. ARCH and GARCH the conditional variance. All shocks were modeled as WN. However, there are models who assume that shocks behaves like a Brownian Motion, mostly in econometry. A Brownian motion or W(t) has four defining properties: it starts at zero, W(0) = 0; it has continuous paths; has independent increments; it has stationary Gaussian increments, which means that $W(t)-W(s)\sim\mathrm{N}(0,t-s),\quad t>s$. W(t) is not stationary, $\mathrm{Var}(W(t)) = t$.  

In 1900, Bachelier modeled the price of a trade asset based on random changes driven by a Brownian Motion. In 1959, Osborne studied Brownian behaviour in log price changes. Samuelson, in 1965, used the GBM (Geometric Brownian motion) model for a positive asset price. In 1973, the Black-Scholes model, authored by Fischer Black and Myron Scholes, explains the theoretical price (premium) of an European call option on a non-dividend-paying stock whose price follows a GBM. In 1998, by Fabienne Comte and Eric Renault, the landmark of stochastic volatility modelling, the FSV (Fractional Stochastic Volatility) model was presented. It modeled the changing, unobservered instanteneous volatility of an asset's returns, with long memory behaviour by choosing a Hurst component (H)$\ > 1/2$. In 2014, with the paper being released in 2018, Gatheral, Jaisson and Rosenbaum presented the RFSV (Rough Fractional Stochastic Volatility) model by choosing the same architecture but using $H < 1/2$. 

\subsubsection{ML Models}
In Machine Learning, the main idea is to create a model who is able to learn data patterns by using a usually larger number of parameters compared to statistical models and a bigger volume of data at once. It's possible the case where the model learns representation already found by statisticians as well as unknown relationships between features. The idea was simple: an outcome is simply a weighted sums over features, we need a way to create those combinations and find the optimal weights to give the expected outcome. The backpropagation algorithm was described by Werbos (1974) and became widely known as the mechanism for neural-network learning (the process of finding the best weights) by Rumelhart et al. in 1986. Several implementations of building the combinations were proposed. 

In 1987, Lapedes and Farber used a Multi-Layer Perceptron (MLP), called as feedforward neural network as well, to take a fixed set of past observations as inputs and learning a non-linear forecasting representation, by using layers of linear combinations transformed by a non-linear activation function. Followed by Waibel et al. in 1989 where a Temporal CNN (Convolutional Neural Network) where convolutions were applied across sucessive time positions to detect patterns, mainly used in speech sequences. 

In 1990, RNN (Recurrent Neural Network) was used to process observations by carrying a hidden state forward and possessing a weigh for the hidden state and another for the inputs. The current state is retrieved by a non-linear combination of inputs and weights. This is similar to an autoregressive model. RNNs had a major issue: they were prone to vanishing gradients, the process where gradients in the backpropagation get zero value or infinitely small values, impeding the weights adjustment from that iteration. Whenever a sequence started to become extremely long, the gradients were becoming smaller and smaller, not being able to reach the weights with the information learnt about the recent steps. In 1997, Hochreiter and Schmidhuber modified the RNN into the LSTM (Long Short Term Memory) model, where they added gates inside the RNN to help retain information over longer spans and built the hidden state from the gates state and the hidden state. They were able to effectively mitigate the issue of longer sequences causing gradients to vanish. 







\subsection{Motivations}
% Explain the problem that motivates the work.
% Example: Unfortunately, there are regions on Earth where there is a shortage of food.

\subsection{Contributions}
% Describe the contributions that address the gaps identified in the motivations.
% Example: We present a method to generate food from thin air.

\subsection{Related Work}
% Describe the state of the art, including relevant algorithms and tools.
% For every paper or tool, explain what this work provides that is not already available.
% When possible, list similar existing tools and explain how this work differs from them.

\subsection{Outline}
% Give an outline of the thesis structure using one short sentence per section.

\section{Background}
% Present all background knowledge required to understand the work.
Unconditional variance measure how the erros vary overall: $\mathrm{Var}(\varepsilon_t)=\sigma^2=\mathrm{E}[(\varepsilon_t-\mathrm{E}[\varepsilon_t])^2]$. For the sample unconditional variance: $s^2=\frac{1}{n-1}\sum_{t=0}^{n-1}(\varepsilon_t-\bar{\varepsilon})^2 $, where $\bar{\varepsilon}$ is the sample mean, equal to $\frac{1}{n}\sum_{t=0}^{n-1}\varepsilon_t$.
\section{Methods}
% Describe the system design and methodology. Do not include source code here.

\section{Implementation}
% Describe how the system was implemented.
% Use pseudocode and/or informal or UML-style diagrams where useful.

\section{Experimental Results}
% Describe the experimental results. This is the most important section.

\subsection{Goals}
% Describe the objectives of the experimental activity, typically:
% - demonstrate correctness;
% - evaluate computational performance, such as CPU/GPU time and RAM usage;
% - evaluate effectiveness, such as expected gain and its standard deviation.

\subsection{Setting}
% Describe the experimental setting: hardware, software, datasets, splits,
% hyperparameters, reproducibility controls, and evaluation metrics.

\subsection{Case Studies}
% Describe the case studies used in the experiments.

\subsection{Correctness}
% Describe experiments that demonstrate the correctness of the implementation.

\subsection{Computational Evaluation}
% Evaluate computational performance, such as training time and CPU, GPU,
% RAM, and VRAM usage.

\subsection{Technical Evaluation}
% Evaluate the effectiveness of the proposed approach.
% Describe the experiments and summarize the results using tables and graphs.

\section{Conclusions}
% Summarize the work and outline possible future developments.

\begin{thebibliography}{99}
% Add one \bibitem entry for every paper or book cited in the main text.
\end{thebibliography}

\end{document}
